\documentclass[11pt]{amsart}

\usepackage[T1]{fontenc}
\usepackage{lmodern}
\usepackage{microtype}
\usepackage{mathtools,amssymb,amsthm}
\usepackage[hidelinks]{hyperref}

\hypersetup{
  pdftitle={Non-monotonicity in k at n = 3: a census contradicting a k-monotonicity
            clause transcribed from D'Angeli and Rodaro}
}

\newtheorem{theorem}{Theorem}
\newtheorem{proposition}[theorem]{Proposition}
\newtheorem{lemma}[theorem]{Lemma}
\newtheorem{corollary}[theorem]{Corollary}
\theoremstyle{remark}
\newtheorem*{remark}{Remark}

\newcommand{\SCC}{\mathrm{SCC}}
\newcommand{\TS}{\mathrm{TS}}

\title[Non-monotonicity in $k$ at $n=3$]
{Non-Monotonicity in $k$ at $n=3$: A Census Contradicting a $k$-Monotonicity
Clause Transcribed from D'Angeli and Rodaro}

\date{}

\subjclass[2020]{05C80, 68Q45, 05C20, 60C05}
\keywords{totally simple graph, lumpable partition, $k$-out digraph, random digraph,
synchronizing automaton, counterexample, exhaustive census}

\begin{document}

\sloppy

\begin{abstract}
For a uniformly random $k$-out digraph $G$ on $n$ labelled vertices --- each vertex
choosing its $k$ out-edge heads independently and uniformly \emph{with replacement} ---
let
$P(n,k)=\Pr\bigl(G \text{ totally simple}\mid G \text{ strongly connected}\bigr)$.
A three-clause conjecture of D'Angeli and Rodaro, as we transcribe it by hand from the
e-print (Section~\ref{sec:locator}), asserts among other things that for each fixed $n$ the
map $k\mapsto P(n,k)$ is nondecreasing on $k\ge2$.  An exhaustive census at $n=3$
contradicts that clause.  Writing $B(n,k)$ for the total weight of the strongly connected states
and $A(n,k)$ for the total weight of those that are also totally simple,
\[
  A(3,2)=114,\quad B(3,2)=296,\qquad A(3,3)=5132,\quad B(3,3)=13754,
\]
so $P(3,2)=\tfrac{57}{148}>\tfrac{2566}{6877}=P(3,3)$, the inequality being decided by
$114\cdot13754=1\,567\,956>1\,519\,072=5132\cdot296$.  All four integers have closed forms
in $k$ that a reader can evaluate by hand, so no computation is needed to check the
inequality.  Two caveats bound what this settles.  The clause, its quantifiers and the
probability model are transcribed by hand from the cited e-print, of which we can supply
neither an immutable copy nor a hash; and the drop, about $1.2$ percentage points,
disappears under the rival uniform measure on labelled $k$-out multigraphs, so
Section~\ref{sec:model} records what pins the with-replacement reading.  Nothing here bears
on the two \emph{asymptotic} clauses of the same conjecture, which remain open.  The
denominators are not new: $B(n,2)$ is a published $1965$/$1971$ automaton census, and
$B(3,3)=13754$ is on record as well.
\end{abstract}

\maketitle

\section{The clause, as transcribed, and its locator}
\label{sec:locator}

The source is D. D'Angeli and E. Rodaro, \emph{On totally synchronizing graphs},
arXiv:2607.17335 \cite{DR}, version~v1, submitted 19 July 2026.  All line references are to
its \LaTeX{} e-print source rather than to the compiled PDF: the file
\texttt{totally\_main.tex}, $143\,403$ bytes, $1977$ lines.  In the section ``Totally
simple graphs'' the source states a three-clause conjecture, tex lines $1457$--$1471$; it
is unnumbered in the source, carries no \verb|\cite|, and is attributed to nobody there.

We can supply neither an immutable copy of that e-print nor a cryptographic hash of it, and
we have not confirmed that the arXiv identifier resolves for a third party.  Everything this
note attributes to the source --- the clause quoted below, its ambient quantifiers, the
definitions of Section~\ref{sec:model} and the probability model --- is hand transcription
from the file named above, and it is not checked by the accompanying program
(Section~\ref{sec:verify}).  A reader who cannot obtain the file cannot check that part.
What is proved here without reference to the source is the arithmetic of
Theorem~\ref{thm:main}: $P(3,2)>P(3,3)$ in the model of Section~\ref{sec:model}.

The clause at issue is the second one, tex line $1463$, transcribed verbatim:

\begin{verbatim}
Moreover, for fixed $n$ the function $k \longmapsto
\Pr(\text{$G$ is totally simple}\mid \mathrm{SCC})$
is nondecreasing.
\end{verbatim}

\noindent
The conjecture opens ``Fix $k\ge2$'', so on this transcription the clause asserts
\begin{equation}\label{eq:clause}
  P(n,k)\;\le\;P(n,k+1)
  \qquad\text{for every $n\ge1$ and every $k\ge2$,}
\end{equation}
where $P(n,k)=\Pr(G\text{ totally simple}\mid\SCC)$.  \emph{The clause places no lower
bound on $n$.}  The first clause of the same conjecture (tex line $1461$) asserts
$P(n,k)\to1$ as $n\to\infty$ for fixed $k$, and the third (tex lines $1465$--$1470$)
asserts an expansion $P(n,k)=1-\exp(-c_kn+\beta_k)+o(1)$ with $c_k$ increasing in $k$.
Both of those are asymptotic; nothing below bears on them.

\section{The model and the predicate}
\label{sec:model}

\subsection*{The measure}
The source fixes the model in the paragraph at tex lines $1445$--$1455$: ``in our model
each vertex chooses independently $k$ outgoing edges with replacement, and the head of
each edge is uniform in $V$''.  So, with $V=\{0,\dots,n-1\}$, the probability space is
$V^{kn}$ with the uniform measure.  Only the out-multiset of each vertex matters for what
follows, so we work with the \emph{state} $c=(c_{vw})_{v,w\in V}$, where $c_{vw}$ is the
number of edges $v\to w$ and $\sum_wc_{vw}=k$ for every $v$; loops and parallel edges are
allowed.  The number of ordered choices realising a state is its weight
\begin{equation}\label{eq:weight}
  W(c)=\prod_{v\in V}\frac{k!}{\prod_{w\in V}c_{vw}!},
  \qquad\text{and}\qquad
  \sum_{c}W(c)=n^{kn}.
\end{equation}
Write $a_{vw}:=c_{vw}$ when the matrix notation is more convenient.

\subsection*{The predicate}
For a partition $\mathcal C$ of $V$ the source calls $\mathcal C$ \emph{lumpable} (tex
$339$--$343$) if for every pair of blocks $C_i,C_j$ and all $u,u'\in C_i$ one has
$|\delta_u(C_j)|=|\delta_{u'}(C_j)|$, where $\delta_u(C_j)$ is the set of edges from $u$
into $C_j$.  A congruence different from the discrete relation $\Delta_G$ and the
universal relation $\nabla_G$ is called \emph{nontrivial} (tex $355$), and an automaton is
\emph{simple} when its only congruences are the discrete and the universal one (tex
$240$).  Chaining that with the definition of \emph{totally simple} (tex $1357$--$1359$),
with the source's lemma identifying lumpable partitions with congruences (tex
$1375$--$1377$) and with its theorem relating lumpability to simplicity (tex
$1386$--$1388$) --- and with the source's own
restatement at tex $1510$, ``total simplicity is equivalent to the absence of nontrivial
lumpable partitions'' --- gives the form we use:
\begin{equation}\label{eq:TS}
\begin{gathered}
  G \text{ is \emph{totally simple}}\\
  \iff\quad
  \text{no partition of $V$ other than $\Delta_G$ and $\nabla_G$ is lumpable.}
\end{gathered}
\end{equation}
Both $\Delta_G$ and $\nabla_G$ are always lumpable ($\nabla_G$ because
$|\delta_u(V)|=k$ for every $u$; $\Delta_G$ vacuously), so excluding both is forced by the
source rather than chosen by us, and the predicate is well posed.

\subsection*{The three quantities}
Set
\begin{equation}\label{eq:ABP}
  B(n,k)=\!\!\sum_{c\ \text{str.\ conn.}}\!\!W(c),
  \qquad
  A(n,k)=\!\!\sum_{\substack{c\ \text{str.\ conn.}\\ c\ \text{totally simple}}}\!\!W(c),
\end{equation}
\begin{equation}\label{eq:P}
  P(n,k)=\frac{A(n,k)}{B(n,k)}.
\end{equation}
Then $P(n,k)=\Pr(G\text{ totally simple}\mid\SCC)$ in the source's model.

\begin{remark}[the model is load-bearing]\label{rem:model}
``Uniformly random $k$-out graph'' also admits a rival reading: the uniform measure on
labelled $k$-out multigraphs, i.e.\ weight $1$ per state instead of \eqref{eq:weight}.
Under that rival reading the counterexample below \emph{does not exist}: writing $P_u$ for
the rival quantity, $P_u(3,2)=\tfrac{34}{65}<\tfrac{133}{243}=P_u(3,3)$, by
$34\cdot486=16\,524<17\,290=266\cdot65$ on the unreduced rival state counts, and every
rival $k$-step we computed is nondecreasing.  Since the drop found below is only about
$1.2$ percentage points, the choice of model is the difference between a counterexample and
none.  Two features of the source pin the ordered reading used here.  First, tex $1449$ says ``with replacement'' in words.
Second, the source's own displayed identity at tex $1452$,
$\Pr(\delta_V(v)=\emptyset)=(1-1/n)^{kn}$ for the event that no edge lands on $v$, is true
in the ordered model and false in the rival one: at $n=3,k=2$ it gives
$(2/3)^6=64/729$, whereas the rival measure gives $(3/6)^3=1/8$.  The residual risk in
this note is therefore a \emph{reading} risk, not an arithmetic one.
\end{remark}

\section{The counterexample at $n=3$}
\label{sec:refutation}

\begin{theorem}\label{thm:main}
In the model of Section~\ref{sec:model},
\[
  A(3,2)=114,\qquad B(3,2)=296,\qquad A(3,3)=5132,\qquad B(3,3)=13754,
\]
so that
\begin{align*}
  P(3,2)&=\frac{114}{296}=\frac{57}{148}=0.385135\ldots,\\
  P(3,3)&=\frac{5132}{13754}=\frac{2566}{6877}=0.373127\ldots
\end{align*}
and therefore $P(3,2)>P(3,3)$.  Consequently \eqref{eq:clause} --- the clause as
transcribed in Section~\ref{sec:locator} from tex line $1463$ of \cite{DR}, read in this
model --- is false.
\end{theorem}

Nothing in the comparison needs decimals.  Cross-multiplying the two unreduced fractions,
\begin{equation}\label{eq:cross}
  114\cdot13754=1\,567\,956
  \;>\;
  1\,519\,072=5132\cdot296,
\end{equation}
and the gap is
\begin{equation}\label{eq:drop}
  P(3,2)-P(3,3)=\frac{57}{148}-\frac{2566}{6877}=\frac{12221}{1017796}=0.0120073\ldots,
\end{equation}
about $1.2$ percentage points.  Here $n=3$ is a fixed $n$ and both $k=2$ and $k=3$ satisfy
the conjecture's own ``Fix $k\ge2$'', so the counterexample lies inside the clause's own
quantifiers.  The value $k=1$ is \emph{not} used: $P(3,1)=1>P(3,2)$ is also a decrease, but
$k=1$ is outside the source's range.

The rest of this section proves Theorem~\ref{thm:main} in closed form, so that a reader
needs no program.  Throughout, $n=3$ and indices are read modulo $3$.

\subsection*{Step 1: total simplicity at $n=3$ is three column inequalities}

\begin{proposition}\label{prop:columns}
A $3$-vertex $k$-out graph with edge multiplicities $a_{ij}$ is totally simple if and only
if in each column the two off-diagonal entries differ:
\[
  a_{10}\neq a_{20},
  \qquad
  a_{01}\neq a_{21},
  \qquad
  a_{02}\neq a_{12}.
\]
\end{proposition}

\begin{proof}
By \eqref{eq:TS} the partitions to be excluded are those other than $\Delta$ and $\nabla$,
and the only such partitions of a $3$-set are the three of shape $2+1$.  Fix
$\mathcal C=\{\{i,i'\},\{j\}\}$.  The singleton block imposes nothing.  For the block
$\{i,i'\}$, lumpability requires $|\delta_i(\{j\})|=|\delta_{i'}(\{j\})|$, i.e.\
$a_{ij}=a_{i'j}$, and $|\delta_i(\{i,i'\})|=|\delta_{i'}(\{i,i'\})|$; since each row sums
to $k$, the latter reads $k-a_{ij}=k-a_{i'j}$ and is equivalent to the former.  So
$\mathcal C$ is lumpable exactly when the two off-diagonal entries of column $j$ agree, and
total simplicity is the conjunction of the three negations.
\end{proof}

\subsection*{Step 2: the denominator $B(3,k)$}

\begin{proposition}\label{prop:B}
$B(3,k)=(3^k-1)^3-3\,(2^k-1)^2(3^k-1)$; in particular
\begin{align*}
  B(3,2)&=8^3-3\cdot3^2\cdot8=512-216=296,\\
  B(3,3)&=26^3-3\cdot7^2\cdot26=17\,576-3822=13\,754,
\end{align*}
and $B(3,4)=80^3-3\cdot15^2\cdot80=512\,000-54\,000=458\,000$.
\end{proposition}

\begin{proof}
For a vertex $v$ let $S_v=\{w\neq v: a_{vw}>0\}$.  A $3$-vertex state fails to be strongly
connected if and only if some $S_v$ is empty (an \emph{all-loop} vertex) or there is a
\emph{closed pair} $\{u,v\}$, meaning $S_u=\{v\}$ and $S_v=\{u\}$.  Indeed, if every $S_v$
is nonempty and the state is not strongly connected, take a sink component $C$ of the
condensation: every vertex of $C$ has its out-support inside $C$, so $|C|=1$ is impossible
and $|C|=3$ would be strong connectivity, leaving $|C|=2$, which is a closed pair.

The ordered weight of the vertices $v$ with $S_v\neq\emptyset$ is $3^k-1$ each (all $3^k$
head choices, minus the single all-loop choice), giving $(3^k-1)^3$ for the first term.
For a fixed closed pair $\{u,v\}$, vertex $u$ must send all $k$ heads into $\{u,v\}$ with
at least one to $v$, of weight $2^k-1$, and likewise for $v$; the third vertex is
unrestricted apart from $S_w\neq\emptyset$, of weight $3^k-1$.  The three closed pairs are
pairwise incompatible, since two of them share a vertex, which would then be forced to be
all-loops and hence to have $S=\emptyset$.  Subtracting gives the formula.
\end{proof}

\subsection*{Step 3: the numerator $A(3,k)$ as a matrix trace}

Give vertex $i$ the ordered pair $(s_i,t_i)=(a_{i,i+1},a_{i,i+2})$, so that
$a_{ii}=k-s_i-t_i$ and the ordered weight of vertex $i$ is the trinomial
\[
  W_k[s][t]=\frac{k!}{s!\,t!\,(k-s-t)!}\quad(s+t\le k),
  \qquad W_k[s][t]=0\quad(s+t>k),
\]
a symmetric $(k+1)\times(k+1)$ matrix.  In these coordinates $a_{10}=t_1$, $a_{20}=s_2$,
$a_{01}=s_0$, $a_{21}=t_2$, $a_{02}=t_0$, $a_{12}=s_1$, so
Proposition~\ref{prop:columns} becomes the cyclic system
\begin{equation}\label{eq:cycle}
  t_0\neq s_1,
  \qquad
  t_1\neq s_2,
  \qquad
  t_2\neq s_0 .
\end{equation}
Let $D=J-I$ be the all-ones-minus-identity matrix of size $k+1$ and put $f(t)=\sum_{q\ge1}W_k[t][q]$.

\begin{proposition}\label{prop:A}
$\displaystyle A(3,k)=\operatorname{tr}\bigl((W_kD)^3\bigr)-3\Bigl[\bigl(\textstyle\sum_tf(t)\bigr)^2-\sum_tf(t)^2\Bigr]$.
\end{proposition}

\begin{proof}
Since $(W_kD)[s][s']=\sum_{t\neq s'}W_k[s][t]$, expanding the trace gives
$\operatorname{tr}((W_kD)^3)=\sum\prod_{i}W_k[s_i][t_i]$ over all
$(s_i,t_i)_{i=0,1,2}$ obeying \eqref{eq:cycle}: that is exactly the total ordered weight of
the totally simple states, strong connectivity not yet imposed.

It remains to subtract the totally simple states that are not strongly connected.  If a
state is totally simple then, by the sink-component argument in the proof of
Proposition~\ref{prop:B} together with Proposition~\ref{prop:columns}, it has no closed
pair: a closed pair $\{u,v\}$ forces $a_{uw}=a_{vw}=0$ in the column of the third vertex
$w$.  Hence a totally simple state fails strong connectivity precisely when some vertex is
all-loops, and at most one vertex can be, since two all-loop vertices $i,i'$ would give
$a_{iw}=a_{i'w}=0$.  Say vertex $0$ is all-loops, i.e.\ $s_0=t_0=0$, of weight
$W_k[0][0]=1$.  Then \eqref{eq:cycle} reads $s_1\ge1$, $t_2\ge1$ and $t_1\neq s_2$, while
the requirement that vertices $1$ and $2$ are not all-loops is automatic.  Summing over
$s_1\ge1$ at fixed $t_1$ contributes $\sum_{s\ge1}W_k[s][t_1]=f(t_1)$ by symmetry of
$W_k$, and summing over $t_2\ge1$ at fixed $s_2$ contributes $f(s_2)$; the surviving
constraint is $t_1\neq s_2$.  So the count is
$\sum_{t_1\neq s_2}f(t_1)f(s_2)=(\sum f)^2-\sum f^2$, and there are three choices of the
all-loop vertex.
\end{proof}

\subsection*{Step 4: the arithmetic, in full}

For $k=2$ and $k=3$ the matrices are, with rows indexed by $s$ and columns by $t$,
\[
  W_2=\begin{pmatrix}1&2&1\\2&2&0\\1&0&0\end{pmatrix},
  \qquad
  W_3=\begin{pmatrix}1&3&3&1\\3&6&3&0\\3&3&0&0\\1&0&0&0\end{pmatrix},
\]
with row sums $4,4,1$ summing to $3^2=9$ and $8,12,6,1$ summing to $3^3=27$.  Their
traces and corrections are
\begin{align*}
  \operatorname{tr}\bigl((W_2D)^3\bigr)&=150,
  & f&=(3,2,0),
  & 3\bigl[5^2-13\bigr]&=36,\\
  \operatorname{tr}\bigl((W_3D)^3\bigr)&=5798,
  & f&=(7,9,3,0),
  & 3\bigl[19^2-139\bigr]&=666,
\end{align*}
so that $A(3,2)=150-36=114$ and $A(3,3)=5798-666=5132$.
Together with Proposition~\ref{prop:B} this proves Theorem~\ref{thm:main}: all four
integers are obtained from a $3\times3$ and a $4\times4$ integer matrix and two
subtractions, and \eqref{eq:cross} concludes.  The same route at $k=4$ gives
$\operatorname{tr}((W_4D)^3)=189\,518$, $f=(15,28,18,4,0)$, correction
$3[65^2-1349]=8628$, hence $A(3,4)=180\,890$ and
\begin{equation}\label{eq:isolated}
  P(3,4)=\frac{180\,890}{458\,000}=\frac{18089}{45800}=0.394956\ldots>P(3,3),
\end{equation}
so monotonicity resumes at the very next step.  The corrections $36$, $666$ and $8628$ are
exactly the weights of the states that are totally simple and \emph{not} strongly
connected.

Propositions~\ref{prop:B} and~\ref{prop:A} give $A(3,k)$ and $B(3,k)$ in closed form for
every $k$.  Evaluating them and comparing consecutive values by exact cross-multiplication
is a finite computation, and the program of Section~\ref{sec:verify} carried it out for
$2\le k\le30$: in that range $k=2\to3$ is the only step along $n=3$ at which
\eqref{eq:clause} fails, and $P(3,k)<P(3,k+1)$ strictly for $3\le k\le29$.  This is a
verified computation rather than a consequence of the closed forms alone, and it says
nothing about $k>30$.

\section{Total simplicity does not imply strong connectivity}
\label{sec:secondary}

This is secondary to the counterexample; it is recorded because the numerator computation
depends on it.

\begin{proposition}\label{prop:1455}
Some totally simple $3$-vertex $2$-out graphs are not strongly connected.  Hence the
sentence transcribed from tex line $1455$ of \cite{DR}, ``if a $k$-out digraph $G$ is not
strongly connected, then it is automatically nontrivially lumpable'', and the assertion of
the source's remark at tex $1516$--$1521$ that every totally simple $k$-out graph is
strongly connected, are both false.
\end{proposition}

\begin{proof}
The $2$-out graph $0\to\{0,0\}$, $1\to\{0,2\}$, $2\to\{1,1\}$ on $\{0,1,2\}$, written as
out-multisets, has $a_{10}=1\neq0=a_{20}$, $a_{01}=0\neq2=a_{21}$ and
$a_{02}=0\neq1=a_{12}$, so it is totally simple by Proposition~\ref{prop:columns}; yet
vertex $0$ is all loops and reaches nothing, so it is not strongly connected.  The remark's
argument takes a sink strongly connected component $C$ and forms ``the equivalence relation
whose only non-singleton class is $C$''; when $|C|=1$ that relation is $\Delta_G$, which the
source itself calls trivial at tex $355$, and the graph just exhibited realises that gap.
\end{proof}

Consequently $P(n,k)$ is \emph{not} $\Pr(\TS)/\Pr(\SCC)$: the joint weight $A(n,k)$ must be
computed with the strong connectivity test applied, as in \eqref{eq:ABP}.

\section{Status: what is settled, and what is not}
\label{sec:status}

The clause \eqref{eq:clause} is false as transcribed: one fixed $n$ and one $k$-step inside
its own quantifiers suffice, and Theorem~\ref{thm:main} exhibits them with closed-form
integers.  Two things bound that.  The clause, its quantifiers and the identification of the
model are hand transcription from an e-print we cannot pin (Section~\ref{sec:locator}); and
the failure is in the ordered/with-replacement model only --- the dip \eqref{eq:drop}
vanishes under the rival reading of Remark~\ref{rem:model}, which records the rival numbers
and the two features of the source that pin the reading used.

The two \emph{asymptotic} clauses of the same conjecture --- $P(n,k)\to1$ as $n\to\infty$
(tex $1461$) and the expansion $1-\exp(-c_kn+\beta_k)+o(1)$ with $c_k$ increasing (tex
$1465$--$1470$) --- are not addressed here, and no finite census can address them.  They
remain open.  The large-$n$ intent of the conjecture survives everything computed here: the
source's own sampled table at $n=8,9,10$ (tex $1484$--$1506$) is monotone in $k$ at each of
those three values of $n$, and $n=3$ lies five below the smallest $n$ it samples.  Replacing
``for fixed $n$'' by ``for fixed $n\ge4$'' is consistent with every cell we know of, but
that is not a theorem: it is quantified over all $n\ge4$ and all $k\ge2$ and supported by
finitely many cells.  Nor do we claim that $k=2\to3$ is the only violation along $n=3$
beyond $k=30$.

The denominators are not new.  The numbers $B(n,2)$ are the classical census of labelled
strongly connected $n$-state $2$-input automata, OEIS \texttt{A027834} \cite{OEIS}, with
$B(3,2)=296$ and $B(4,2)=20\,958$ the third and fourth terms; the sequence goes back to
Harrison \cite{Harrison} and Liskovets \cite{Liskovets}.  Neither is $B(3,3)=13\,754$ new.
An editorial note in OEIS \texttt{A006691} \cite{OEISb}, quoting Robinson \cite{Robinson}
and Liskovets \cite{Liskovets}, records \texttt{A006692} ($49$, $6877$, $1854545,\dots$) as
a normalised count of labelled strongly connected automata with three inputs, the $p$-state
value being $(p-1)!$ times its $(p-1)$st term; at $p=3$ that is $2!\cdot6877=13\,754$.  On
that reading $B(n,k)$ is on record for every $k$, not only for $k=2$.  We have not verified
that normalisation independently, and the program checks only the two terms $296$ and
$20\,958$ against its own census.

Adjacent work settles nothing about the numerator or about $k$-monotonicity.  Harrison's
census of strongly connected sequential machines \cite{HarrisonSCC} and the subsequent
literature on automorphism groups and quotients of strongly connected automata count or
classify such automata without the lumpability refinement, and the sources just cited give
$B(n,k)$ but not $A(n,k)$; none of them addresses the behaviour of $P(n,k)$ in $k$.  We
claim no priority for the lumpability-refined numerator $A(n,k)$: searches of
\texttt{A027834}'s host database along both the $k$-axis ($114,5132,180890$) and the
$n$-axis ($1,9,114,6720$) return no sequence, which is a failed search and not evidence of
novelty.

\section{Reproduction and verification}
\label{sec:verify}

The accompanying program \texttt{verify.py} (Python~3.9+, standard library only)
independently re-derives the census integers of Theorem~\ref{thm:main} and the other
census values printed here, within the limits recorded below.  It enumerates the state space
exhaustively at $n=3$ for $k=1,2,3,4$, at $n=4$ for $k=2$ and at $n=2$ for $k=2,3$,
deciding lumpability from the definition over \emph{all} set partitions with $\Delta$ and
$\nabla$ removed; it re-runs $n=3$, $k=2,3$ a second time as a direct enumeration over the
$3^{kn}$ raw ordered edge tuples with weight $1$ per tuple and no multinomial weight
anywhere; it checks the two independent closed forms of Propositions~\ref{prop:B}
and~\ref{prop:A} against the census and against each other for $2\le k\le30$; it verifies
the total-weight identity $n^{kn}$ of \eqref{eq:weight} in every cell and that $\Delta$ and
$\nabla$ are lumpable on every state; and it decides every inequality by
exact integer or \texttt{Fraction} arithmetic.  Its recorded run reports
\texttt{VERDICT: ALL 56 CHECKS PASS}.

The limits of that program are recorded in its own output and repeated here.  It does not
recompute any cell with $n\ge5$, nor $n=4$ with $k\ge3$, so the repair discussed in
Section~\ref{sec:status} is consistent with, not verified by, it; and it does not recompute
the source's Monte Carlo table at $n=8,9,10$, nor any attribution in
Section~\ref{sec:status} --- of those it checks only that its own census reproduces the two
hand-transcribed terms $296$ and $20\,958$, and it does not touch the
\texttt{A006691}/\texttt{A006692} normalisation.  Above all, it does not read \cite{DR}: the
quoted clause, the tex line numbers, the definitions of Section~\ref{sec:model} and the
identification of the probability model are transcribed from the e-print by hand, and the
program checks arithmetic in the model as transcribed.  Its output banner speaks of a
counterexample to the clause and calls the clause refuted; that reading of the source is the
part of this note a reader must check against \cite{DR} directly, and it is the residual
risk recorded in Remark~\ref{rem:model} and in Section~\ref{sec:locator}.

\begin{thebibliography}{9}

\bibitem{DR}
D.~D'Angeli and E.~Rodaro,
\emph{On totally synchronizing graphs},
arXiv:2607.17335v1, 19 July 2026.
Abstract page \href{https://arxiv.org/abs/2607.17335v1}{arxiv.org/abs/2607.17335v1};
e-print source \href{https://arxiv.org/e-print/2607.17335v1}{arxiv.org/e-print/2607.17335v1}.
Line numbers quoted in this note refer to the \LaTeX{} source file
\texttt{totally\_main.tex} of that e-print.  We supply no copy or hash of that file.

\bibitem{Harrison}
M.~A. Harrison,
\emph{A census of finite automata},
Canad. J. Math. \textbf{17} (1965), 100--113.

\bibitem{HarrisonSCC}
M.~A. Harrison,
\emph{Enumeration of strongly connected sequential machines},
Inform. and Control \textbf{8} (1965),
\href{https://doi.org/10.1016/S0019-9958(65)90316-5}{doi:10.1016/S0019-9958(65)90316-5}.

\bibitem{Robinson}
R.~W. Robinson,
\emph{Counting strongly connected finite automata},
in: Graph Theory with Applications to Algorithms and Computer Science
(Kalamazoo, 1984), Wiley, New York, 1985.

\bibitem{Liskovets}
V.~A. Liskovets,
\emph{The number of connected initial automata},
Vesti Akad. Nauk Belarus. SSR, Ser. Fiz.-Mat. Navuk \textbf{3} (1971), 26--30.

\bibitem{OEIS}
OEIS Foundation Inc.,
\emph{Sequence \texttt{A027834}: number of labeled strongly connected $n$-state $2$-input
automata},
The On-Line Encyclopedia of Integer Sequences,
\href{https://oeis.org/A027834}{oeis.org/A027834}.

\bibitem{OEISb}
OEIS Foundation Inc.,
\emph{Sequences \texttt{A006691} (with the editorial note of P.~Hadjicostas, February
$2021$) and \texttt{A006692}: connected $n$-state finite automata with $3$ inputs},
The On-Line Encyclopedia of Integer Sequences,
\href{https://oeis.org/A006691}{oeis.org/A006691},
\href{https://oeis.org/A006692}{oeis.org/A006692}.

\end{thebibliography}

\end{document}
